\chapter{Abelian Varieties}\label{abelian-varieties}
\chapterstatus{complete}

\section{Introduction}\label{abelian-varieties:section-introduction}

A complex abelian variety is a complex torus $\CC^g/\Lambda$ that embeds in projective space. Tori are the simplest compact K\"ahler manifolds, and their
cohomology is completely explicit: exterior powers of $H^1$. For this reason abelian varieties are the main testing ground of the Hodge conjecture. They
include the known cases (divisors, Lefschetz classes), the first instances of exceptional Hodge classes (Weil classes), and Deligne's theorem that Hodge classes
on abelian varieties are absolutely Hodge (Part~\ref{part-hodge-theory}). This chapter proves Riemann's criterion for a torus to be algebraic, describes line
bundles, polarisations, the dual and the Jacobian, and discusses endomorphisms and complex multiplication. References: \cite{birkenhake-lange},
\cite{mumford-abelian}, \cite[Chapter 2, Section 6]{griffiths-harris}.

\noindent\textbf{Prerequisites.} Chapter~\ref{curves} (Algebraic Curves and Riemann Surfaces), Chapter~\ref{divisors} (Divisors and Line Bundles),
Chapter~\ref{hodge-decomposition} (The Hodge Decomposition) and Chapter~\ref{algebraic-groups} (Linear Algebraic Groups).

Throughout, $V$ is a complex vector space of dimension $g$, $\Lambda \subseteq V$ a lattice (a discrete subgroup of rank $2g$ spanning $V$ over $\RR$), and $T = V/\Lambda$ the complex
torus, with its structure of compact complex Lie group (group law written additively). $\bar V^\vee$ denotes the space of conjugate-linear forms $V \to \CC$.

\section{Complex tori and their cohomology}\label{abelian-varieties:section-tori}

\begin{proposition}\label{abelian-varieties:proposition-torus-cohomology}
\begin{enumerate}
\item $H_1(T; \ZZ) = \Lambda$, and cup product gives ring isomorphisms $H^k(T; \ZZ) \cong \Lambda^k\Hom(\Lambda, \ZZ)$, the alternating $k$-forms on $\Lambda$.
\item The Hodge decomposition is $H^1(T; \CC) = \Hom_\RR(\Lambda \otimes \RR, \CC) = V^\vee \oplus \bar V^\vee$, with $H^{1,0} = V^\vee$ (spanned by the $dz_j$) and
$H^{0,1} = \bar V^\vee$ (spanned by the $d\bar z_j$), and $H^{p,q}(T) = \Lambda^pV^\vee \otimes \Lambda^q\bar V^\vee$.
\end{enumerate}
\end{proposition}

\begin{proof}
(1) $T$ is diffeomorphic to $(S^1)^{2g}$ via a $\ZZ$-basis of $\Lambda$, and $\RR^{2g} \to T$ is the universal covering, so $\pi_1(T) = H_1(T; \ZZ) = \Lambda$. The ring structure follows from
the K\"unneth formula (Theorem~\ref{singular-cohomology:theorem-kunneth}) and $H^\bullet(S^1) = \Lambda^\bullet\ZZ$. (2) With a flat metric the harmonic forms are the constant-coefficient
forms (Example~\ref{dolbeault:example-hodge-numbers}); a constant $1$-form is a real-linear form on $V$, and it splits uniquely into complex-linear and conjugate-linear parts.
\end{proof}

\begin{definition}\label{abelian-varieties:definition-abelian-variety}
A complex torus is an \emph{abelian variety} if it is projective. A homomorphism of tori is a holomorphic group homomorphism, and an \emph{isogeny} is a surjective homomorphism
with finite kernel. More generally, over any field an \emph{abelian variety} is a complete connected group variety; such a group is automatically commutative and projective
\cite[Chapter II]{mumford-abelian}.
\end{definition}

\begin{lemma}\label{abelian-varieties:lemma-homomorphisms}
Every holomorphic map $f \colon T \to T'$ of complex tori with $f(0) = 0$ is a homomorphism, induced by a complex-linear map $F \colon V \to V'$ with $F(\Lambda) \subseteq \Lambda'$.
\end{lemma}

\begin{proof}
Lift $f$ to $F \colon V \to V'$ with $F(0) = 0$. For $\lambda \in \Lambda$, $F(v + \lambda) - F(v) \in \Lambda'$ depends continuously on $v$, hence is constant. So the partial derivatives
of $F$ are $\Lambda$-periodic holomorphic functions, so they descend to holomorphic functions on the compact connected manifold $T$, which are constant by the maximum principle
(Proposition~\ref{several-complex-variables:proposition-identity-maximum}). So $F$ is linear, and $F(\Lambda) \subseteq \Lambda'$ because $F(\lambda) = F(0 + \lambda) - F(0) \in \Lambda'$.
\end{proof}

\section{Line bundles and the Riemann conditions}\label{abelian-varieties:section-riemann}

\begin{proposition}\label{abelian-varieties:proposition-neron-severi}
Identify $H^2(T; \ZZ)$ with the alternating forms $E \colon \Lambda \times \Lambda \to \ZZ$ (Proposition~\ref{abelian-varieties:proposition-torus-cohomology}). A class $E$ lies in
$H^{1,1}(T)$ if and only if $E(iu, iv) = E(u, v)$ for all $u, v \in V$ (extending $E$ $\RR$-bilinearly), if and only if $E = \operatorname{Im}H$ for a (unique) Hermitian form $H$ on $V$,
namely $H(u, v) = E(iu, v) + iE(u, v)$. Every such $E$ is $c_1(L)$ for a holomorphic line bundle $L$, so
\[
\mathrm{NS}(T) := \operatorname{im}(c_1) = \{H \text{ Hermitian on } V : \operatorname{Im}H(\Lambda, \Lambda) \subseteq \ZZ\}.
\]
\end{proposition}

\begin{proof}
A real alternating form $E$ on $V$ is a constant real $2$-form. Its $(2,0)$ and $(0,2)$ components vanish exactly when $E$ is of type $(1,1)$ for the complex structure $i$,
that is, $E(iu, iv) = E(u, v)$: the $(2,0) + (0,2)$ part is $\frac12(E - E(i\cdot, i\cdot))$. Given this, $H(u, v) = E(iu, v) + iE(u, v)$ is Hermitian:
$H(iu, v) = E(-u, v) + iE(iu, v) = i(E(iu, v) + iE(u, v)) = iH(u, v)$, and $\overline{H(v, u)} = E(iv, u) - iE(v, u) = E(iu, v) + iE(u, v)$, using $E(iv, u) = -E(u, iv) = -E(iu, i^2v) = E(iu, v)$.
Conversely $\operatorname{Im}$ of a Hermitian form satisfies the condition. The last statement is Lemma~\ref{kodaira:lemma-integral-11}.
\end{proof}

\begin{theorem}[Riemann]\label{abelian-varieties:theorem-riemann-conditions}
A complex torus $T = V/\Lambda$ is an abelian variety if and only if there is a positive definite Hermitian form $H$ on $V$ with $\operatorname{Im}H(\Lambda, \Lambda) \subseteq \ZZ$.
Such an $H$ (or its alternating form $E = \operatorname{Im}H$) is a \emph{polarisation}; it is the first Chern class of an ample line bundle.
\end{theorem}

\begin{proof}
For a constant real $2$-form $\omega$ put $E_\omega = -\omega$ and $H_\omega(u, v) = E_\omega(iu, v) + iE_\omega(u, v)$. If $\omega$ is of type $(1,1)$, then $H_\omega$ is Hermitian
(Proposition~\ref{abelian-varieties:proposition-neron-severi}) with $\operatorname{Re}H_\omega(u, u) = -\omega(iu, u) = \omega(u, iu)$. So $\omega$ is a positive $(1,1)$-form, that is, the
fundamental form of a Hermitian metric (Proposition~\ref{kahler:proposition-fundamental-form-coordinates}), if and only if $H_\omega$ is positive definite. Conversely every Hermitian $H$ is
$H_\omega$ for $\omega = -\operatorname{Im}H$.

Suppose $H$ is positive definite with $\operatorname{Im}H$ integral on $\Lambda$. Then $\omega = -\operatorname{Im}H$ is a positive constant $(1,1)$-form, closed because its coefficients
are constant, and its class is integral. So $T$ is a Hodge manifold (Definition~\ref{kodaira:definition-hodge-manifold}) and projective by
Corollary~\ref{kodaira:corollary-projective-criterion}.

Conversely, let $T$ be projective, with an ample line bundle $L$ and a metric of positive curvature form $\omega_h$ (Theorem~\ref{gaga:theorem-projective-kahler}). Average over translations:
$\bar\omega = \int_Tt_x^*\omega_h\,dx$ for the normalised Haar measure (Theorem~\ref{lie-groups:theorem-haar}). Each $t_x^*\omega_h$ is positive and cohomologous to $\omega_h$, because
translations are homotopic to the identity, so $\bar\omega$ is a positive translation-invariant closed $(1,1)$-form in the integral class $c_1(L)$. A translation-invariant form has constant
coefficients. So $H_{\bar\omega}$ is positive definite, and $\operatorname{Im}H_{\bar\omega} = -\bar\omega$ is integral on $\Lambda$, being (up to sign) the integral class $c_1(L)$ in
$H^2(T; \ZZ) = \Lambda^2\Hom(\Lambda, \ZZ)$.
\end{proof}

\begin{remark}\label{abelian-varieties:remark-sign-convention}
The proof shows that ample line bundles correspond to positive definite $H$ once a constant form $\omega$ is matched with the alternating form $-\omega$; this is the convention of
\cite{birkenhake-lange}, where $c_1(L) = \operatorname{Im}H$. With the opposite matching one gets negative definite forms instead; Riemann's condition is insensitive to this, since
$-H$ is positive definite if $H$ is negative definite.
\end{remark}

\begin{proposition}\label{abelian-varieties:proposition-period-matrix}
Let $E$ be a polarisation. There is a basis $\lambda_1, \ldots, \lambda_g, \mu_1, \ldots, \mu_g$ of $\Lambda$ with
$E(\lambda_j, \mu_k) = d_j\delta_{jk}$ and $E(\lambda_j, \lambda_k) = E(\mu_j, \mu_k) = 0$, where $d_1 \mid d_2 \mid \cdots \mid d_g$ are positive integers, the \emph{type} of the polarisation.
It is \emph{principal} if all $d_j = 1$. Then $\frac{1}{d_1}\mu_1, \ldots, \frac{1}{d_g}\mu_g$ is a $\CC$-basis of $V$, and in it $\Lambda = Z\ZZ^g + D\ZZ^g$ with $D = \mathrm{diag}(d_j)$
and $Z$ a symmetric complex $g \times g$ matrix with $\operatorname{Im}Z$ positive definite. Conversely every such $Z$ in the \emph{Siegel upper half space} $\mathfrak{H}_g$ defines a
polarised abelian variety of type $D$.
\end{proposition}

\begin{proof}
The symplectic basis is the elementary divisor theorem for alternating forms over $\ZZ$ (Frobenius), proved like the Smith normal form (Theorem~\ref{rings:theorem-smith}); see
\cite{birkenhake-lange}. Since $H$ is positive definite, the real span of the $\mu_j$ is $E$-isotropic, and $H$ is real and positive on it, so the $\mu_j$ are $\CC$-linearly
independent: a complex relation $\sum c_j\mu_j = 0$ with $c = a + ib$ gives $\sum a_j\mu_j = -i\sum b_j\mu_j$, and pairing with $H$ forces $a = b = 0$. Writing
$\lambda_j = \sum_kZ_{kj}\mu_k/d_k$, the conditions that the $\lambda$'s span an isotropic subspace and that $H$ is positive definite translate into $Z = Z^T$ and
$\operatorname{Im}Z > 0$ by a direct computation \cite{birkenhake-lange}. The converse is the same computation read backwards together with
Theorem~\ref{abelian-varieties:theorem-riemann-conditions}.
\end{proof}

\begin{example}\label{abelian-varieties:example-dimension-one}
For $g = 1$, $\mathfrak{H}_1$ is the upper half plane, and $\Lambda = \ZZ\tau + \ZZ$ with $\operatorname{Im}\tau > 0$. Every one-dimensional torus is an abelian variety, as Riemann's
condition only asks $\operatorname{Im}\tau > 0$ (compare Example~\ref{lefschetz-theorems:example-curves-riemann}). The alternating form $E(\tau, 1) = 1$ gives
$H(u, v) = u\bar v/\operatorname{Im}\tau$.
\end{example}

\begin{counterexample}\label{abelian-varieties:counterexample-non-algebraic-torus}
For $g \geq 2$ the general complex torus is not algebraic. Fix a real basis identifying $V$ with $\RR^{2g}$ and $\Lambda$ with $\ZZ^{2g}$, so that complex tori correspond to complex
structures $J$ on $\RR^{2g}$. For each of the countably many nonzero integral alternating forms $E$, the condition $E(Ju, Jv) = E(u, v)$ defines a proper closed real-analytic subset of the
space of complex structures when $g \geq 2$ \cite{birkenhake-lange}. Outside the countable union of these subsets,
$\mathrm{NS}(T) = 0$ by Proposition~\ref{abelian-varieties:proposition-neron-severi}, so there is no ample line bundle and $T$ is not projective. For such tori the Hodge conjecture
in degree $2$ has no analytic cycles to work with; more generally the analytic version of the conjecture fails for some complex tori (Chapter~\ref{failures}).
\end{counterexample}

\begin{theorem}[Appell--Humbert]\label{abelian-varieties:theorem-appell-humbert}
Holomorphic line bundles on $T$ correspond bijectively to pairs $(H, \alpha)$ with $H$ Hermitian, $E = \operatorname{Im}H$ integral on $\Lambda$, and $\alpha \colon \Lambda \to U(1)$ a
\emph{semicharacter}: $\alpha(\lambda + \mu) = \alpha(\lambda)\alpha(\mu)e^{i\pi E(\lambda, \mu)}$. The bundle $L(H, \alpha)$ is the quotient of $V \times \CC$ by
$\lambda \cdot (v, t) = (v + \lambda, \alpha(\lambda)e^{\pi H(v, \lambda) + \frac{\pi}{2}H(\lambda, \lambda)}t)$, and $c_1(L(H, \alpha)) = E$. Its sections are the
\emph{theta functions}, and if $H$ is positive definite of type $(d_1, \ldots, d_g)$, then $h^0(L) = d_1 \cdots d_g$. Moreover $L^2$ is globally generated and $L^3$ is very ample
(Lefschetz).
\end{theorem}

This is \cite{birkenhake-lange}; see also \cite{mumford-abelian}.

\section{The dual abelian variety}\label{abelian-varieties:section-dual}

\begin{proposition}\label{abelian-varieties:proposition-dual}
$\Pic^0(T) \cong \hat T = \bar V^\vee/\hat\Lambda$, where $\hat\Lambda = \{l \in \bar V^\vee : \operatorname{Im}l(\Lambda) \subseteq \ZZ\}$, a complex torus of dimension $g$, the
\emph{dual torus}. A polarisation $H$ defines a homomorphism $\varphi_H \colon T \to \hat T$, $v \mapsto H(v, \cdot)$, which is an isogeny of degree $(d_1 \cdots d_g)^2$, and an
isomorphism if the polarisation is principal. If $T$ is an abelian variety, so is $\hat T$.
\end{proposition}

\begin{proof}
By the exponential sequence, $\Pic^0(T) = H^1(T, \mathcal{O})/H^1(T; \ZZ)$ (Theorem~\ref{holomorphic-bundles:theorem-exponential}), and $H^1(T, \mathcal{O}) = H^{0,1} = \bar V^\vee$
(Proposition~\ref{abelian-varieties:proposition-torus-cohomology}). The map $H^1(T; \RR) = \Hom_\RR(V, \RR) \to \bar V^\vee$ sends a real linear form $a$ to its conjugate-linear part
$l_a$, and $a = 2\operatorname{Re}l_a$. So the image of $H^1(T; \ZZ) = \Hom(\Lambda, \ZZ)$ is $\{l : 2\operatorname{Re}l(\Lambda) \subseteq \ZZ\}$, which the $\CC$-linear automorphism
$l \mapsto 2il$ of $\bar V^\vee$ carries onto $\hat\Lambda$, since $\operatorname{Im}(2il) = 2\operatorname{Re}l$. The map
$v \mapsto H(v, \cdot)$ is complex-linear from $V$ to $\bar V^\vee$ and maps $\Lambda$ into $\hat\Lambda$ because $\operatorname{Im}H(\Lambda, \Lambda) \subseteq \ZZ$. It is injective since
$H$ is nondegenerate, so its kernel on $T$ is $\hat\Lambda'/\Lambda$ with $\hat\Lambda' = \{v : \operatorname{Im}H(v, \Lambda) \subseteq \ZZ\}$, of order $\det E = (d_1 \cdots d_g)^2$ in
a symplectic basis. The inverse of $H$ gives a polarisation on $\hat T$ \cite{birkenhake-lange}.
\end{proof}

\section{Jacobians}\label{abelian-varieties:section-jacobians}

\begin{theorem}\label{abelian-varieties:theorem-jacobian}
Let $C$ be a compact Riemann surface of genus $g$. The Jacobian $J(C) = H^0(C, \Omega^1)^\vee/H_1(C; \ZZ)$ (Definition~\ref{curves:definition-jacobian}) is a principally
polarised abelian variety, the polarisation being the intersection form on $H_1(C; \ZZ)$.
\end{theorem}

\begin{proof}
$J(C)$ is a complex torus of dimension $g$ (Proposition~\ref{curves:proposition-jacobian-torus}). The intersection form $E$ on $H_1(C; \ZZ)$ is unimodular (Poincar\'e duality,
Corollary~\ref{poincare-duality:corollary-cup-pairing}), so a symplectic basis $a_1, \ldots, a_g, b_1, \ldots, b_g$ exists with $E(a_j, b_k) = \delta_{jk}$. Extend $E$ $\RR$-linearly to
$V = H^0(\Omega^1)^\vee \cong H_1(C; \RR)$. Under Poincar\'e duality $H_1(C; \RR) \cong H^1(C; \RR)$ and the real isomorphism $H^1(C; \RR) \to H^{1,0}$, $\alpha \mapsto \alpha^{1,0}$,
the form $E$ becomes $\pm\int_C\alpha \wedge \beta = \pm 2\operatorname{Re}\int_C\alpha^{1,0} \wedge \overline{\beta^{1,0}}$, and the complex structure of $V$ becomes that of $H^{1,0}$ (up to
conjugation). So $E(iu, iv) = E(u, v)$, and the Hermitian form $H(u, v) = E(iu, v) + iE(u, v)$ is, up to sign and a positive factor, $i\int_C\alpha^{1,0} \wedge \overline{\beta^{1,0}}$, which
is definite by the Hodge--Riemann relations (Theorem~\ref{lefschetz-theorems:theorem-hodge-riemann}, Example~\ref{lefschetz-theorems:example-curves-riemann}). Replacing $E$ by $-E$ if
necessary, $H$ is positive definite, and $J(C)$ is an abelian variety by Theorem~\ref{abelian-varieties:theorem-riemann-conditions}. The polarisation is principal as $E$ is unimodular.
\end{proof}

\begin{remark}\label{abelian-varieties:remark-torelli-schottky}
By Torelli's theorem the principally polarised Jacobian determines $C$. The Jacobians form a subvariety of dimension $3g - 3$ of the moduli space $\mathfrak{H}_g/\mathrm{Sp}_{2g}(\ZZ)$ of
principally polarised abelian varieties, which has dimension $g(g + 1)/2$. For $g \leq 3$ the dimensions agree and the general principally polarised abelian variety is a Jacobian;
for $g \geq 4$ it is not, and characterising Jacobians is the \emph{Schottky problem} \cite{birkenhake-lange}.
\end{remark}

\section{Endomorphisms and complex multiplication}\label{abelian-varieties:section-endomorphisms}

\begin{proposition}\label{abelian-varieties:proposition-elliptic-endomorphisms}
Let $E = \CC/(\ZZ\tau + \ZZ)$. Then $\End(E)$ is $\ZZ$, or an order in the imaginary quadratic field $\QQ(\tau)$; the latter happens if and only if $\tau$ is quadratic over $\QQ$.
\end{proposition}

\begin{proof}
By Lemma~\ref{abelian-varieties:lemma-homomorphisms} an endomorphism is multiplication by some $\alpha \in \CC$ with $\alpha\Lambda \subseteq \Lambda$. Then $\alpha = a + b\tau$ and
$\alpha\tau = c + d\tau$ with integers $a, b, c, d$. If $b \neq 0$, eliminating $\alpha$ gives $b\tau^2 + (a - d)\tau - c = 0$, so $\tau$ is quadratic, imaginary as $\operatorname{Im}\tau > 0$,
and $\alpha \in \QQ(\tau)$ is integral (it is an eigenvalue of an integer matrix). If $b = 0$ then $\alpha = a \in \ZZ$. Conversely, if $a\tau^2 + b\tau + c = 0$ with integers $a \neq 0$,
$b$, $c$, then $\alpha = a\tau$ satisfies $\alpha \cdot 1 = a\tau \in \Lambda$ and $\alpha\tau = -b\tau - c \in \Lambda$, so multiplication by $\alpha \notin \ZZ$ is an endomorphism.
\end{proof}

\begin{definition}\label{abelian-varieties:definition-cm}
An abelian variety $A$ of dimension $g$ has \emph{complex multiplication} (CM) if $\End(A) \otimes \QQ$ contains a commutative semisimple subalgebra of dimension $2g$ over $\QQ$,
for instance a CM field $F$ (a totally imaginary quadratic extension of a totally real field) with $[F : \QQ] = 2g$.
\end{definition}

\begin{theorem}\label{abelian-varieties:theorem-endomorphism-algebra}
Let $A$ be an abelian variety over $\CC$.
\begin{enumerate}
\item (Poincar\'e reducibility) $A$ is isogenous to a product $A_1^{n_1} \times \cdots \times A_r^{n_r}$ with the $A_i$ simple and pairwise non-isogenous, and
$\End(A) \otimes \QQ \cong \prod_iM_{n_i}(D_i)$ with $D_i = \End(A_i) \otimes \QQ$ division algebras.
\item A polarisation defines the \emph{Rosati involution} $f \mapsto f^\dagger$ on $\End(A) \otimes \QQ$, and $\operatorname{tr}(ff^\dagger) > 0$ for $f \neq 0$.
\item (Shimura--Taniyama) If $A$ has CM by a CM field $F$ of degree $2g$, then $A \cong \CC^g/\Phi(\mathfrak{a})$ for a fractional ideal $\mathfrak{a}$ of $F$ and a \emph{CM type}
$\Phi = \{\phi_1, \ldots, \phi_g\}$, a set of embeddings $F \to \CC$ containing exactly one of each pair of complex conjugate embeddings.
\end{enumerate}
\end{theorem}

This is \cite{birkenhake-lange} and \cite{mumford-abelian}. The endomorphism algebras $D_i$ are classified by
Albert into four types. For the Hodge conjecture, CM abelian varieties are the case where the Mumford--Tate group is a torus (Chapter~\ref{mumford-tate}), and where exceptional
Hodge classes first appear (Chapter~\ref{known-cases}).

\begin{example}\label{abelian-varieties:example-cm}
\begin{enumerate}
\item $E_i = \CC/\ZZ[i]$ has $\End(E_i) = \ZZ[i]$ and $E_\rho = \CC/\ZZ[\rho]$, $\rho = e^{2\pi i/3}$, has $\End(E_\rho) = \ZZ[\rho]$.
\item For $E$ without CM, $\End(E \times E) \otimes \QQ = M_2(\QQ)$, and $\mathrm{NS}(E \times E) \otimes \QQ$ has rank $3$; for $E$ with CM it has rank $4$. This is
Example~\ref{divisors:example-picard-numbers}(3).
\item The Jacobian of the curve $y^2 = x^5 - 1$ has CM by $\QQ(\zeta_5)$, a CM field of degree $4 = 2g$.
\end{enumerate}
\end{example}

\section{Exercises}\label{abelian-varieties:section-exercises}

\begin{exercise}\label{abelian-varieties:exercise-ns-product}
Let $E_1$, $E_2$ be elliptic curves. Show that $\mathrm{NS}(E_1 \times E_2) \cong \ZZ^2 \oplus \Hom(E_1, E_2)$, and deduce that its rank is $2$, $3$ or $4$.
\end{exercise}

\begin{solution}
Write $\Lambda = \Lambda_1 \oplus \Lambda_2$. An integral alternating form of type $(1,1)$ on $\Lambda$ restricts to such forms on $\Lambda_1$ and $\Lambda_2$, each a multiple of the generator
($\ZZ^2$), and its mixed part $E(\lambda_1, \lambda_2)$ is a pairing of type $(1,1)$ between $\Lambda_1$ and $\Lambda_2$, which via $E_1$'s principal polarisation is the same as a
$\ZZ$-linear map $\Lambda_1 \to \Lambda_2$ commuting with $i$ after identification, that is, a homomorphism $E_1 \to E_2$ (Lemma~\ref{abelian-varieties:lemma-homomorphisms}). By
Proposition~\ref{abelian-varieties:proposition-elliptic-endomorphisms} and its analogue for homomorphisms, $\Hom(E_1, E_2)$ has rank $0$, $1$ or $2$.
\end{solution}

\begin{exercise}\label{abelian-varieties:exercise-hodge-classes-torus}
Show that the Hodge classes in $H^2(T; \QQ)$ of a complex torus form the space $\mathrm{NS}(T) \otimes \QQ$, and that for $T = E^g$, $E$ without CM, this space has dimension
$g(g + 1)/2$.
\end{exercise}

\begin{solution}
The first statement is Proposition~\ref{abelian-varieties:proposition-neron-severi}. For $E^g$, $\mathrm{NS} \otimes \QQ$ is the space of $g \times g$ Hermitian matrices with entries in
$\End(E) \otimes \QQ = \QQ$ that are Hermitian for the Rosati involution, that is, symmetric rational matrices, of dimension $g(g + 1)/2$ \cite{birkenhake-lange}.
\end{solution}

\begin{exercise}\label{abelian-varieties:exercise-isogeny-dual}
Show that multiplication by $n$ on $T$ is an isogeny of degree $n^{2g}$ with kernel $\frac1n\Lambda/\Lambda \cong (\ZZ/n)^{2g}$.
\end{exercise}

\begin{solution}
The map is induced by $v \mapsto nv$ on $V$, surjective, with kernel $\{v : nv \in \Lambda\}/\Lambda = \frac1n\Lambda/\Lambda$, which has $n^{2g}$ elements since $\Lambda \cong \ZZ^{2g}$.
\end{solution}
